% paper/sections/06_collocate.tex
\section{コロケート格子と圧力速度連成問題}
\label{sec:collocate}

\subsection{コロケート格子の利点と課題}
本研究では，すべての物理量（速度 $\bu$, 圧力 $p$, レベルセット $\phi, \psi$）を同一の格子点（セル中心）に配置する\textbf{コロケート格子}（collocated grid）を採用する．この配置は，複雑な形状の境界を扱う際に実装が単純であり，高次精度スキーム（例えば本稿のCCD法）をすべての変数に統一的に適用できるという大きな利点を持つ．

しかし，コロケート格子には特有の課題が存在する．それは，\textbf{圧力と速度の数値的な分離（decoupling）}に起因する，非物理的な振動解（チェッカーボードパターン）が発生しうることである．

\begin{tcolorbox}[defbox, title={なぜチェッカーボード問題が起きるのか？：1次元の例}]
1次元の運動量方程式（$x$方向）を考えよう．セル $i$ における圧力勾配を単純な中心差分で近似すると，
\[ \left(\pd{p}{x}\right)_i \approx \frac{p_{i+1} - p_{i-1}}{2\Delta x} \]
この式からわかるように，セル $i$ の運動量は，隣接するセル $i+1$ と $i-1$ の圧力しか参照しない．セル $i$ 自身の圧力 $p_i$ は計算に寄与しないのである．

ここで，圧力場が `... +1, -1, +1, -1, +1, -1 ...` のような波長 $2\Delta x$ の激しい振動（チェッカーボード）を持つ場合を想像してみよう．
このとき，どのセル $i$ で圧力勾配を計算しても，
\[ \frac{p_{i+1} - p_{i-1}}{2\Delta x} = \frac{(+1) - (+1)}{2\Delta x} = 0 \quad \text{または} \quad \frac{(-1) - (-1)}{2\Delta x} = 0 \]
となり，圧力勾配は常にゼロと評価されてしまう．つまり，\textbf{速度場はこのような高周波の圧力振動を全く「感知」できない}．
結果として，本来ならば滑らかな圧力場が得られるべき問題で，物理的に意味のないギザギザの圧力解が誤差として残り続けてしまう．これが圧力速度デカップリング問題の本質である．
\end{tcolorbox}

\subsection{Rhie--Chow 補正による問題の解決}
この問題を解決するため，1983年に Rhie と Chow によって提案された補間法が広く用いられている．
その本質は，セル中心の離散化運動量方程式を適切に補間して，セル界面（face）の速度を定義することにある．
これにより，すべての格子点における圧力が速度の計算に寄与するようになり，デカップリングが抑制される．

\subsubsection{補正された界面速度の定義}
本稿では，Chung (2002) に基づく変密度場に対応した形式を用いる．
東側フェイス $e$ における $x$ 方向速度 $u_e$ は，隣接するセル中心 $P$ と $E$ の値を用いて以下のように定義される．
\begin{equation}
 u_e = \overline{u_e} - d_e \left[ \left(\pd{p}{x}\right)_e - \overline{\left(\pd{p}{x}\right)_e} \right]
 \label{eq:rc-concept}
\end{equation}
ここで，
\begin{itemize}
    \item $\overline{u_e}$: セル中心 $P, E$ の速度を線形補間した速度．$\overline{u_e} = (u_P+u_E)/2$．
    \item $(\partial p/\partial x)_e$: フェイス $e$ での圧力勾配．単純な中心差分 $\frac{p_E-p_P}{\Delta x}$ で計算．
    \item $\overline{(\partial p/\partial x)_e}$: セル中心 $P, E$ で評価された圧力勾配を線形補間したもの．
    $\overline{(\partial p/\partial x)_e} = \frac{1}{2}\left[ \left(\frac{\partial p}{\partial x}\right)_P + \left(\frac{\partial p}{\partial x}\right)_E \right]$
    \item $d_e$: 補間係数．運動量方程式から導かれ，時間刻み幅 $\Delta t$ と密度 $\rho$ に依存する．
    $d_e \approx (\Delta t / \rho)_e$.
\end{itemize}

\paragraph{補正項の働き}
式 \eqref{eq:rc-concept} の第2項が Rhie-Chow 補正の核心である．
この項は「\textbf{フェイス上で直接計算した圧力勾配}」と「\textbf{セル中心の圧力勾配を補間したもの}」との\textbf{差}に比例する．
\begin{itemize}
    \item \textbf{滑らかな圧力場の場合}：
    圧力場が滑らかであれば，2つの圧力勾配の評価方法はほぼ同じ値を与えるため，補正項は自然にゼロに近くなる．
    \item \textbf{チェッカーボード圧力場の場合}：
    圧力場がギザギザの場合，フェイスでの勾配 $(p_E-p_P)/\Delta x$ は大きな値を持つが，セル中心での勾配の平均 $\overline{(\partial p/\partial x)_e}$ はゼロに近くなる．
    この差が大きな補正量となり，非物理的な振動を効果的に\textbf{減衰（damping）}させる働きをする．
\end{itemize}

\subsubsection{本研究における実装形式と変密度場への拡張}
最終的に，圧力 Poisson 方程式との整合性を厳密に保ち，かつ，気液二相流のような密度が不連続に変化する場を安定して扱うため，本研究では以下の実装形式を用いる．
\begin{equation}
 u_f = \bar{u}_f - \Delta t\left( \frac{1}{\rho} \right)_f^{\text{avg}} \left[ (\bnabla p)_f - \overline{(\bnabla p)}_f \right]
 \label{eq:rc-face}
\end{equation}
ここで，最も重要な変更点は，補正係数 $d_f \approx \Delta t/\rho_f$ の計算に，密度の逆数の\textbf{算術平均}（arithmetic mean）を用いている点である．

\begin{tcolorbox}[defbox, title={なぜ変密度場で「密度の逆数の調和平均」を用いるのか？}]
空気（$\rho_g \approx 1$）と水（$\rho_l \approx 1000$）のように密度が$1000:1$にもなる気液界面を考えよう．
この界面がちょうどセルフェイス $f$ 上にあるとする．
このとき，フェイスにおける圧力勾配流束 $\frac{1}{\rho}\nabla p$ をどう評価すれば物理的に正しくなるだろうか？

\paragraph{算術平均（単純な線形補間）の破綻}
もし密度 $\rho$ を単純に線形補間（算術平均）すると，$\rho_f = (\rho_l+\rho_g)/2 \approx 500$ となる．
この値は液相でも気相でもなく，物理的実態を表していない．
さらに深刻なのは，圧力補正後の速度発散から圧力Poisson方程式を導出する際に，運動量保存の観点から問題が生じることである．
定常状態では，フェイスを通過する質量流束 $\rho u$ と圧力勾配流束 $\frac{1}{\rho}\nabla p$ は連続でなければならない．
算術平均ではこの流束の連続性が保証されず，特に高密度比の場合，界面で非物理的な圧力振動や速度誤差を引き起こすことが知られている．

\paragraph{調和平均による流束の連続性の保証}
この問題を解決するのが調和平均である．2点 $P, N$ における値を持つ係数 $k$ の調和平均は $\frac{2k_P k_N}{k_P+k_N}$ で定義される．
「密度の逆数」$1/\rho$ の調和平均をフェイス $f$ で求めると，
\[ \left( \frac{1}{\rho} \right)_f^{\text{harm}} = \frac{2 (1/\rho_P)(1/\rho_N)}{(1/\rho_P)+(1/\rho_N)} = \frac{2}{\rho_P+\rho_N} \cdot \frac{\rho_P\rho_N}{1} \]
ではなく，正しくはフェイスにおける「伝熱係数」と同様に，フェイスを挟む2セルの抵抗の和として考える．
これにより，フェイスにおける $(1/\rho)$ の値は以下のように定義されるべきである．
\[ \left(\frac{1}{\rho}\right)_f \approx \frac{1}{2}\left( \frac{1}{\rho_P} + \frac{1}{\rho_N} \right) \]
これは密度の逆数 $1/\rho$ の\textbf{算術平均}に相当する．
（注：しばしば文献で「調和平均」と引用されるのは密度 $\rho$ そのものの調和平均であり，それは $1/\rho$ の算術平均と等価ではないが，同様の物理的要請から導かれる．本稿の実装は $1/\rho$ の算術平均を用いている．）

この定義を用いると，高密度側（水, $\rho_l$）の寄与が支配的だった算術平均とは対照的に，低密度側（空気, $\rho_g$）の大きな値 $1/\rho_g$ が強く影響する．
これにより，圧力勾配に対する速度応答が両相で適切にスケーリングされ，界面をまたぐ流束の連続性が離散的にも保証される．
結果として，高密度比の界面においても安定した圧力速度連成が達成されるのである．
\end{tcolorbox}
